\section*{الفصل \ref{ch:b1:rifts-ridges-margins} --- الشقوق والحيود والهوامش السلبية}

\begin{solution}{exo:b1:rifts-ridges-margins:1}
(أ) يغطي السدّ في المسقط $\qty{60000}{m}\times\qty{8}{m} = \qty{4.8e5}{m^2}$، فيكون ارتفاعه المتوسط $\qty{2.5e9}{m^3}/\qty{4.8e5}{m^2} =
\qty{5200}{m}$، أي نحو \qty{5}{km}: فالسدّ قطع جزءًا كبيرًا من \omterm{def:b1:layered-earth:layers}{القشرة}. (ب) $\qty{8}{m}/\qty{0.018}{m.yr^{-1}} = 444$ سنة، أي نحو أربعة قرون.
\end{solution}

\begin{solution}{exo:b1:rifts-ridges-margins:2}
(أ) \omterm{def:b1:cross-sections:throw}{الإزاحة الأفقية} هي \omterm{def:b1:cross-sections:throw}{الإزاحة الرأسية} مقسومةً على ظلّ \omterm{def:b1:geological-maps:attitude}{الانحدار}:
$3.0/\tan\ang{60} = \qty{1.73}{km}$ و$0.8/\tan\ang{70} = \qty{0.29}{km}$؛
فيكون امتداد المقطع \qty{2.02}{km}. (ب) قبل الصدوع كان عرض المقطع $25 - 2.02 = \qty{22.98}{km}$، أي نحو \qty{23}{km}: فقد استطال بنحو \qty{9}{\percent}.
\end{solution}

\begin{solution}{exo:b1:rifts-ridges-margins:3}
فائقة البطء (أقل من \qty{20}{mm.yr^{-1}}): حيد غاكيل (12)، ووسط البحر الأحمر (15)، وحيد ريكيانيس (19)، وهذا الأخير قريب من الحدّ. والبطيئة (من 20 إلى \qty{55}{mm.yr^{-1}}): حيد كارلزبرغ الغربي (23)، والأطلسي الجنوبي (33). والسريعة (أكثر من \qty{75}{mm.yr^{-1}}): حيد شرق المحيط الهادئ (143). وينبغي أن يكون للحيود فائقة البطء والبطيئة وادٍ محوري بين صدوع عادية، ولحيد شرق المحيط الهادئ قمّة محورية.
\end{solution}

\begin{solution}{exo:b1:rifts-ridges-margins:4}
(أ) $\qty{3.0}{km^2}\times\qty{7.1}{km} = \qty{21.3}{km^3}$، أي نحو \qty{21}{km^3} في السنة. (ب) $\qty{3.028e8}{km^2}/\qty{3.0}{km^2.yr^{-1}} =
\qty{1.01e8}{yr}$، أي \qty{101}{Myr}. وأقدم قاع المحيط، نحو \qty{190}{Myr}، أقدم من ذلك لأن القاع لا يُفنى بترتيب الأعمار: فالاندساس يستهلك قاعًا من كل عمر، ويبقى بعض القاع القديم حيث لم يصل إليه خندق، وقد تغيّر معدّل الإنتاج عبر الزمن.
\end{solution}

\begin{solution}{exo:b1:rifts-ridges-margins:5}
(أ) $2500 + 350\times2 = \qty{3200}{m}$، و$2500 + 350\times6 = \qty{4600}{m}$،
$2500 + 350\times8 = \qty{5300}{m}$. (ب) $t = \bigl((d - 2500)/350\bigr)^2$:
$(2100/350)^2 = \qty{36}{Myr}$ و$(2800/350)^2 = \qty{64}{Myr}$. (ج) ينخفض القاع \qty{1400}{m} بين 4 و\qty{36}{Myr}، لكن \qty{700}{m} فقط بين 36 و\qty{64}{Myr}، في مدّة زمنية مماثلة.
\end{solution}

\begin{solution}{exo:b1:rifts-ridges-margins:6}
معامل التمدّد الفعّال هو
$\num{3.2e-5}\times3300/2300 = \qty{4.591e-5}{K^{-1}}$،
$\sqrt{\kappa\times\qty{1}{Myr}} = \sqrt{\num{e-6}\times\num{3.156e13}} =
\qty{5618}{m}$. ثم $C = (2/\sqrt\pi)\times\num{4.591e-5}\times1350\times5618 =
1.1284\times\num{4.591e-5}\times1350\times5618 = \qty{393}{m.Myr^{-1/2}}$. وبالنسبة إلى \omterm{def:b1:layered-earth:layers}{وشاح} أبرد بمقدار \qty{250}{K}، يتناسب $C$ مع $T_m - T_0$، فيكون:
$393\times1100/1350 = \qty{320}{m.Myr^{-1/2}}$. ويهبط قاع المحيط الأبرد أبطأ، ويتبع الثابت الأصغر. والرقم \num{350} المركّب في الفصل ليس من هذه الحالة: فهو يقع ضمن عدم اليقين في الثوابت.
\end{solution}

\begin{solution}{exo:b1:rifts-ridges-margins:7}
(أ) $\qty{e-6}{cal.cm^{-2}.s^{-1}} = \num{4.184e-6}\times\num{e4} =
\qty{4.184e-2}{W.m^{-2}} = \qty{41.84}{mW.m^{-2}}$، فيكون الثابت $11.3\times41.84 = \qty{473}{mW.m^{-2}.Myr^{1/2}}$، أي أدنى بنحو \qty{3.5}{\percent} من القيمة 490 في \cref{ch:b1:heat-budget}: فالتركيبان يتفقان. (ب) $q = 473$، و$473/\sqrt{10} = 150$ و$473/\sqrt{150} = \qty{38.6}{mW.m^{-2}}$. ويواصل قانون الجذر التربيعي الهبوط، بينما يستوي نموذج الصفيحة عند قيمة ثابتة: فعند \qty{150}{Myr} تكون قيمة الجذر التربيعي أدنى من اللازم، إذ يبلغ قانون العمق عمقًا أكبر مما ينبغي.
\end{solution}

\begin{solution}{exo:b1:rifts-ridges-margins:8}
خُذ القطعة الشمالية على $x = 0$، والجنوبية على $x = -40$ (\unit{km})، والتحويل على $y = 0$ بينهما. (أ) مباشرةً شمال التحويل يقع القاع غرب القطعة الشمالية ويتحرك غربًا؛ ومباشرةً جنوبه يقع شرق القطعة الجنوبية ويتحرك شرقًا. ومن أيّ من الجانبين يتحرك الجانب البعيد نحو اليسار: فالانزلاق من النوع اليساري. (ب) يبدو أن الحيد أُزيح يمينيًا، وجزؤه الجنوبي إلى الغرب: وهذا الاتجاه المعاكس، كما تتنبأ به \cref{prop:b1:rifts-ridges-margins:transform}. (ج) تقع الزلازل فقط على امتداد \qty{40}{km} بين طرفي الحيد، حيث ينتمي الجانبان إلى صفيحتين مختلفتين؛ وخارجه، على امتداد مناطق الكسر، ينتمي الجانبان إلى الصفيحة نفسها ويتحركان معًا.
\end{solution}

\begin{solution}{exo:b1:rifts-ridges-margins:9}
المعدّل النصفي هو \qty{10.6}{mm.yr^{-1}}، أي \qty{10.6}{km.Myr^{-1}}. (أ) فرق العمر هو $350/10.6 = \qty{33.0}{Myr}$. (ب) عند نهاية التحويل يقابل القاعُ البالغ عمره \qty{33.0}{Myr} ($2500 + 350\sqrt{33.0} = \qty{4511}{m}$) قاعًا عمره 0 (\qty{2500}{m}): فرق قدره نحو \qty{2010}{m}. وعلى بعد \qty{100}{km} بعده، تكون الأعمار $450/10.6 = 42.5$ و$100/10.6 = \qty{9.4}{Myr}$، والأعماق 4781 و\qty{3575}{m}: فرق قدره نحو \qty{1210}{m}. وعلى بعد \qty{1000}{km} بعده، تكون الأعمار 127.4 و\qty{94.3}{Myr}، والأعماق 6450 و\qty{5899}{m}: \qty{550}{m}. (هذه الأعمار تتجاوز \qty{80}{Myr}، حيث يبلغ قانون الجذر التربيعي عمقًا زائدًا؛ ويعطي قانون الصفيحة فرقًا أصغر.)
\end{solution}

\begin{solution}{exo:b1:rifts-ridges-margins:10}
(أ) مع مساحات متساوية لكل عمر، يكون العمق المتوسط هو متوسط $d(t)$ على $[0, A]$:
\[
  \frac1A\int_0^A\bigl(d_r + C\sqrt t\bigr)\mathrm dt = d_r + \frac CA\cdot\frac23A^{3/2}
  = d_r + \frac23C\sqrt A.
\]
(ب) من أجل $A = \qty{180}{Myr}$: $2500 + \frac23\times350\times13.416 =
\qty{5630}{m}$؛ وعند \qty{120}{Myr}: $2500 + \frac23\times350\times10.954 =
\qty{5056}{m}$، أي أضحل بمقدار \qty{574}{m}. (ج) لنفترض أن مستوى البحر يرتفع بمقدار $r$ وأن القاع المحمّل يهبط بمقدار $s = r\rho_w/(\rho_m - \rho_w)$. فالماء، الذي لا يتغيّر حجمه، يجب أن يملأ العمق نفسه فوق المساحة نفسها، إذن $r + s = \qty{574}{m}$، أي $r\,\rho_m/(\rho_m - \rho_w) = 574$، و$r = 574\times2300/3300 =
\qty{400}{m}$: مئات الأمتار، بمعدّل الانتشار وحده.
\end{solution}

\begin{solution}{exo:b1:rifts-ridges-margins:11}
(أ) $\sqrt{x - a} = \sqrt x\sqrt{1 - a/x} \approx \sqrt x\,(1 - a/2x)$ من أجل
$a \ll x$، إذن
$\Delta d \approx C\sqrt{x/u}\cdot a/(2x) = Ca/(2\sqrt{ux})$.
(ب) عند $x = \qty{450}{km}$: الفرق الدقيق \qty{1205}{m}، والتقريبي
$350\times350/(2\sqrt{10.6\times450}) = \qty{887}{m}$؛ وعند \qty{1350}{km}:
\qty{550}{m} و\qty{512}{m}؛ وعند \qty{3350}{km}: \qty{334}{m} و
\qty{325}{m}. ونسبة التقريبي إلى الدقيق لا تتعلّق إلا بالمقدار $r = a/x$: $(r/2)/(1 - \sqrt{1 - r})$، وهي 0.9 عند $r = 0.36$. ويكون التقريب جيدًا بنسبة \qty{10}{\percent} وراء $x = 350/0.36 \approx
\qty{970}{km}$، أي نحو ثلاث إزاحات عن الحيد.
\end{solution}

\begin{solution}{exo:b1:rifts-ridges-margins:12}
(أ) $d_2 - d_1 = 3671.1 - 3606.8 = \qty{64.3}{m}$ عند \qty{10}{Myr}؛ $4072.8 - 4065.3 = \qty{7.5}{m}$ عند 20؛ $4415.4 - 4417.0 = \qty{-1.7}{m}$ عند 30؛ $4707.5 - 4713.6 = \qty{-6.1}{m}$ عند 40. (ب) $d_2 - d_1$ موجب عند 20 وسالب عند 30، والقانونان متصلان، فهما متساويان بينهما. التنصيف: عند 25، $d_2 - d_1 = +0.9$؛ وعند 27.5، $-0.7$؛ وعند 26.25، $+0.01$؛ وعند 26.875، $-0.3$؛ وعند 26.56، $-0.2$؛ وعند 26.41، $-0.08$؛ وعند 26.33، $-0.04$: فالجذر بين 26.25 و26.33، عند \qty{26.3}{Myr}، حيث يعطي القانونان كليهما \qty{4294}{m}. (ج) $d_1' = 175/\sqrt{26.3} = \qty{34.1}{m.Myr^{-1}}$ و$d_2' = (3200/62.8)\,e^{-26.3/62.8} = \qty{33.5}{m.Myr^{-1}}$. فللقانونين القيمة نفسها تقريبًا والميل نفسه هناك، فلا يظهر للمنحنى الموصول أي انكسار.
\end{solution}

\begin{solution}{pb:b1:rifts-ridges-margins:1}
\textbf{1.} \omterm{def:b1:cross-sections:throw}{الإزاحة الأفقية} $= $ \omterm{def:b1:cross-sections:throw}{الإزاحة الرأسية}$/\tan$(\omterm{def:b1:geological-maps:attitude}{الانحدار}): F1 \qty{346}{m}، وF2 \qty{770}{m}،
F3 \qty{2014}{m}، وF4 \qty{163}{m}، وF5 \qty{462}{m}، وF6 \qty{1172}{m}. فالصدع F3، ذو أكبر \omterm{def:b1:cross-sections:throw}{إزاحة رأسية} وأفقية، هو صدع الحدّ.

\textbf{2.} \omterm{def:b1:geological-maps:attitude}{الامتداد} هو المجموع، \qty{4927}{m}، أي نحو \qty{4.9}{km}. وكان الشقّ قبل الصدوع عرضه $60 - 4.9 = \qty{55.1}{km}$.

\textbf{3.} العمودي على شقّ يمتد باتجاه \ang{020} يشير إلى \ang{110}؛ والحركة، نحو \ang{105}، تصنع معه \ang{5}. والمركبة العمودية هي $6.0\cos\ang{5} = \qty{5.98}{mm.yr^{-1}}$.

\textbf{4.} $\qty{4.927}{km}/\qty{5.98}{km.Myr^{-1}} = \qty{0.82}{Myr}$.

\textbf{5.} في \qty{3.2}{Myr} تباعدت الصفيحتان عبر الشقّ بمقدار $5.98\times3.2 = \qty{19.1}{km}$. وتحمل الصدوع منه $4.93/19.1 = \qty{26}{\percent}$.

\textbf{6.} الأرباع الثلاثة الأخرى، أي \qty{14.2}{km}، ليست في صدوع المقطع. وقد تستوعبها سدود تتسلّل على امتداد المحور، فتفتح الشقّ دون أن تصدعه، كما في عفار، أو صدوع خارج المقطع. ومع سدود عرض كل منها \qty{8}{m}، يلزم نحو $\qty{14200}{m}/\qty{8}{m} \approx 1780$ سدًّا، واحدًا كل $3.2\times10^6/1780 \approx 1800$ سنة في المتوسط.

\textbf{7.} من C2An: $(21.0 + 20.2)/2.581 = 41.2/2.581 =
\qty{15.96}{mm.yr^{-1}}$. ومن C3n: $(41.0 + 40.2)/5.235 = 81.2/5.235 =
\qty{15.51}{mm.yr^{-1}}$.

\textbf{8.} $(41.2 - 2.2)/2.581 = \qty{15.11}{mm.yr^{-1}}$ و
$(81.2 - 2.2)/5.235 = \qty{15.09}{mm.yr^{-1}}$: بعد التصحيح يتفق المعدّلان في حدود \qty{0.02}{mm.yr^{-1}}. ويكون للتصحيح أثر أكبر في الحقبة الأحدث، لأن مسافتها أقصر.

\textbf{9.} تزيح الإزاحة إلى الخارج حافتي الجانب الواحد بالمقدار نفسه، فتلغى في فرقهما. وعلى الجانبين معًا،
$(81.2 - 41.2)/(5.235 - 2.581) = 40.0/2.654 = \qty{15.07}{mm.yr^{-1}}$.

\textbf{10.} شمال شرق: $(41.0 - 21.0)/2.654 = \qty{7.54}{mm.yr^{-1}}$؛ وجنوب غرب: $(40.2 - 20.2)/2.654 = \qty{7.54}{mm.yr^{-1}}$. فالانتشار متماثل؛ والفرق \qty{0.8}{km} بين الجانبين، وهو نفسه للحقبتين، إزاحة ثابتة، مثل وضع المحور \qty{0.4}{km} بعيدًا عن المركز.

\textbf{11.} $\cos\Delta = \sin\ang{30.9}\sin\ang{18.0} +
\cos\ang{30.9}\cos\ang{18.0}\cos\ang{16.4} = 0.15870 + 0.78287 = 0.94157$، إذن
$\Delta = \ang{19.69}$. ومع $\omega R = 0.403\times(\pi/180)\times6371 =
\qty{44.81}{mm.yr^{-1}}$، يكون المعدّل $44.81\sin\ang{19.69} =
\qty{15.10}{mm.yr^{-1}}$، أي المعدّلان المصحّحان في السؤالين 8 و9.

\textbf{12.} $\cos\Delta = \sin\ang{30.9}\sin\ang{26.0} +
\cos\ang{30.9}\cos\ang{26.0}\cos\ang{11.4} = 0.22512 + 0.75600 = 0.98112$، إذن
$\Delta = \ang{11.15}$، والمعدّل $44.81\sin\ang{11.15} =
\qty{8.7}{mm.yr^{-1}}$.

\textbf{13.} تتطابق التنبؤات، 15.1 و\qty{8.7}{mm.yr^{-1}}، مع المقاس 15 و\qty{9}{mm.yr^{-1}}. فالمعدّل هو $\omega R\sin\Delta$، ويقع القطب شمال غرب البحر الأحمر، في شمال شرق ليبيا: فكلما اتجهنا جنوبًا ابتعدنا عن القطب، وكان الانفتاح أسرع.

\textbf{14.} $15.09\times5.235 = \qty{79}{km}$ من قاع المحيط.

\textbf{15.} $5000/120.6 = \qty{41}{km.Myr^{-1}}$، أي \qty{41}{mm.yr^{-1}} في المتوسط.

\textbf{16.} $5000/15.1 = \qty{331}{Myr}$ بالمعدّل الأوسط؛
$5000/8.7 = \qty{575}{Myr}$ بمعدّل الشمال.

\textbf{17.} لأنه ينفتح حول قطب قريب منه، يتّسع البحر الأحمر كمقصّ: يزداد عرضه متناسبًا مع $\sin\Delta$، ويبقى إسفينًا يضيق نحو القطب. وعند \qty{8.7}{mm.yr^{-1}}، يحتاج الجزء الشمالي إلى $79/8.7 = \qty{9.1}{Myr}$ ليصنع \qty{79}{km} من قاع المحيط التي يملكها الجزء الأوسط الآن.

\textbf{18.} عند \qty{15}{mm.yr^{-1}} يكون محور البحر الأحمر حيدًا فائق البطء؛ ويُتوقّع وادٍ محوري بين صدوع عادية.

\textbf{19.} $2500 + 350\times1 = \qty{2850}{m}$;
$2500 + 350\sqrt{10} = \qty{3607}{m}$; $2500 + 350\sqrt{40} = \qty{4714}{m}$.

\textbf{20.} $t = (1500/350)^2 = \qty{18.4}{Myr}$.

\textbf{21.} في أول \qty{10}{Myr} ينخفض بمقدار $350\sqrt{10} = \qty{1107}{m}$؛
وبين 70 و\qty{80}{Myr}، $350(\sqrt{80} - \sqrt{70}) = 350\times(8.944 -
8.367) = \qty{202}{m}$، أي أقل بنحو خمس مرات في المدّة نفسها.

\textbf{22.} $\mathrm dd/\mathrm dt = C/(2\sqrt t) = 175/\sqrt{80} =
\qty{19.6}{m.Myr^{-1}}$، أي نحو \qty{0.02}{mm} في السنة.

\textbf{23.} $490/\sqrt{80} = \qty{54.8}{mW.m^{-2}}$.

\textbf{24.} قانون الجذر التربيعي: $2500 + 350\sqrt{80} = 2500 + 350\times8.944 =
\qty{5630}{m}$. وقانون الصفيحة: $6400 - 3200\,e^{-80/62.8} = 6400 - 3200\times0.2798
= \qty{5505}{m}$. والفرق \qty{126}{m}. وعند \qty{80}{Myr} ينطبق القانونان بعد: قانون الجذر التربيعي حتى نحو \qty{80}{Myr}، وقانون الصفيحة من نحو \qty{20}{Myr}؛ وبعد ذلك يبلغ قانون الجذر التربيعي عمقًا زائدًا.

\textbf{25.} يعطي قانون الجذر التربيعي \qty{5630}{m}: بأقرب عُشر من الكيلومتر، يقع قاع المحيط الجديد على عمق \textbf{\qty[text-series-to-math=true]{5.6}{km}} عندما يبلغ 80 مليون سنة (ويعطي قانون الصفيحة \qty{5.5}{km}).
\end{solution}
